
\begin{figure*}[t]
    \centering
    \begin{subfigure}[b]{0.95\linewidth}
        \centering
        \includegraphics[width=\linewidth]{figs/vllm+diffusion_performance_bar.pdf}
        \vspace{-2em}
        \caption{Multimodal AI (Text Pipeline): vLLM (high priority) + Diffusion (low priority).}
        \label{fig:eval:vllm+diffusion_mempressure_bar}
        \vspace{0.3em}
    \end{subfigure}
    \begin{subfigure}[b]{0.95\linewidth}
        \centering
        \includegraphics[width=\linewidth]{figs/vllm+diffusion.vl_performance_bar.pdf}
        \vspace{-2em}
        \caption{Multimodal AI (Video Pipeline): vLLM (high priority) + Diffusion (low priority).}
        \label{fig:eval:vllm+diffusion_mempressure_bar}
        \vspace{0.3em}
    \end{subfigure}
    
    \begin{subfigure}[b]{0.95\linewidth}
        \centering
        \includegraphics[width=\linewidth]{figs/vllm+llama_factory_performance_bar.pdf}
        \vspace{-2em}
        \caption{Continual online learning (Text Pipeline): vLLM (high priority) + LlamaFactory (low priority).}
        \label{fig:eval:vllm+llama_factory_mempressure_bar}
    \end{subfigure}
    \begin{subfigure}[b]{0.95\linewidth}
        \centering
        \includegraphics[width=\linewidth]{figs/vllm+llama_factory.vl_performance_bar.pdf}
        \vspace{-2em}
        \caption{Continual online learning (Video Pipeline): vLLM (high priority) + LlamaFactory (low priority).}
        \label{fig:eval:vllm+llama_factory_mempressure_bar}
    \end{subfigure}
    \vspace{-1em}
    \caption{Performance variation when co-running a high priority (HP) task and a low priority (LP) one under memory pressure. \tool achieves the lowest latency that is close to HP running exclusively, and high LP throughput. \yq{tmp updated; Yicheng has latest numbers to update; Ghost/GhostDYN is transparent/signal-based, we should make it more explicit.} 
    %\kyle{I am a bit concerned about GhostDYN, it's more like trading LLM latency for diffusion throughput. I was expecting both of them will become better by leveraging application semantics.}
    }
    \label{fig:eval:mempressure-overall}
    % \vspace{-1em}
\end{figure*}

\mysection{Evaluation}
\label{sec:eval}

Our evaluation aims to answer the following questions:
\squishlist
    \item Does \tool provide robust performance isolation and adaptively reallocate resources in a way that yields better performance than existing approaches? (\S\ref{sec:eval:e2e})
    \item Can \tool flexibly trade-off latency slackness to best-effort throughput and stay on the Pareto frontier? (\S\ref{sec:eval:tradeoff})
    \item What contributes to \tool's superior performance? (\S\ref{sec:eval:ablation})
\squishend

\mysubsection{Setup}

\MyPara{Testbed.}
We conducted experiments on a GCP instance with one NVIDIA A100-40G GPU and 85GB host memory. The server runs Ubuntu 24.04 (Linux 6.14.0), CUDA 12.9, and NVIDIA open GPU kernel modules 575.57.08.

\MyPara{Baselines.}
We compare \tool against state-of-the-art prior work including TGS~\cite{tgs@nsdi23} (transparent memory oversubscription), XSched~\cite{xsched@osdi25} (user-space scheduling), and GPreempt~\cite{gpreempt@atc25} (kernel-level preemption). We also evaluate industry-standard hardware partitioning using MIG~\cite{nvidiaMIG}.

TGS is the only baseline that natively supports memory overcommitment. For others, we enable memory overcommitment by intercepting their CUDA APIs and routing allocation through UVM, following TGS's approach.
Note that TGS builds upon MPS~\cite{nvidiaMPS} and UVM optimizations, serving as a strictly stronger baseline than standard MPS. We have verified that the TGS consistently achieves the same or higher performance than MPS. 

% \MyPara{Models and frameworks.}

\mysubsection{Applications}

% \yq{Are these applications sufficient?}
% \yq{We can have the same frameworks but with different workloads. for example, use vllm to serve different models.}

We constructed two AI pipelines which cover three frameworks, three models with both LLM and diffusion models, training and inference, different priority levels, and different SLO requirements. 
These pipelines reflects emerging trends in multi-modal agentic systems and new continual learning workloads. 
% \hx{Again, say how representative these applications are of the real-world scenarios..}
% \joey{I added a direct statement that is likely to be non-controversial?}

% \hx{Kyle, write something here about the omni-vllm or -sglang to further motivate the setup.}
\MyPara{Multi-modal AI} integrates the capability of LLMs and diffusion models to process and generate contents in both text and other modalities such as images and videos. The rapid shift toward compound AI systems has driven major serving frameworks to introduce ``omni'' variants\textemdash such as vLLM-omni~\cite{vllm-omni} and SGLang-omni~\cite{sglang-omni}\textemdash designed to natively orchestrate text, vision, and audio generation on shared GPU infrastructure. By co-locating a latency-critical LLM with a throughput-oriented diffusion model, we can faithfully replicate the complex resource contention and scheduling dynamics inherent to these evolving omni-modal deployments.

We built a pipeline for this with a Llama-3.1 8B LLM and a stable-diffusion 3.5-medium diffusion model. The pipeline consists of a vLLM instance and a customized image generation instance built with Hugging Face diffusers library~\cite{huggingfaceDiffusers}. Since LLM inference is more latency-critical, it runs at a high priority. Diffusion requires a longer time for high-quality images with a loose or no latency requirement, so we treat it as a low-priority task.
For this pipeline, we used BurstGPT trace~\cite{burstgpt} for LLM and VidProM dataset~\cite{wang2024vidprommillionscalerealpromptgallery} for the diffusion model.

% Frameworks: vLLM + Diffusers
% Models: Llama-3B/8B for vLLM, sd3.5-medium for Diffusers
% Workloads: same trace for inference.

% Limit Diffusers to ensure vLLM always get enough memory and avoid swapping in most times. We also set a fixed limit for vLLM to ensure its latency.

\MyPara{Continual online learning}
is an emerging LLM deployment paradigm that enables a model to keep learning from new data over time during inference~\cite{jayasuriya2025sparcsubspaceawarepromptadaptation,pcl-continual-learning@www25}. In a continual learning pipeline, an LLM is served for both inference and training with periodical weights update.
To this end, we build such a pipeline for Llama-3.2-3B with vLLM~\cite{vllm@sosp23} for inference and Llama-Factory~\cite{zheng2024llamafactory} for LoRA finetuning. vLLM runs at a high priority to ensure low inference latency, while Llama-Factory serves as a best-effort task and runs only when vLLM cannot fully utilize the GPU.
Following prior studies~\cite{blitzscale@osdi25}, we use the BurstGPT~\cite{burstgpt} trace for inference and the Alpaca dataset~\cite{alpaca@github} for finetuning.

%\hx{Where is the third pipeline?}
%\yq{We only have two pipelines now.}


\begin{comment}
\MyPara{Multi-LLM Serving} is gaining more attention recently due to the surge in interest in multi-agent systems. A multi-LLM pipeline typically consists of several different LLMs that work jointly for higher-quality task handling.
We simulate this with Qwen3-8B served by vLLM~\cite{vllm@sosp23} and Llama-3.1-8B with Llama.cpp~\cite{githubGitHubGgmlorgllamacpp}. Here we use different serving frameworks to showcase \tool's generality.
\end{comment}

\mysubsection{End-to-End Performance}
\label{sec:eval:e2e}

We first compare \tool against all baselines in both memory-constrained and memory-sufficient settings.

\mysubsubsection{Performance Isolation}

We run both pipelines in a memory-constrained configuration to test whether \tool can preserve high-priority performance while maintaining useful low-priority throughput.
Results are shown in Figure~\ref{fig:eval:mempressure-overall}.
When co-locating vLLM (HP) and Diffusion (LP) (Figure~\ref{fig:eval:vllm+diffusion_mempressure_bar}), \tool achieves the lowest vLLM latency and the highest diffusion throughput.
Compared to the best-performing baseline, \tool reduces vLLM median and P99 TTFT and ITL by 59$\times$, 5.1$\times$, 1.2$\times$, and 2.8$\times$, 
respectively, while keeping all four metrics within 15\% of the ideal latency achieved by vLLM running exclusively. At the same time, \tool delivers the highest diffusion throughput among all systems, demonstrating that strong isolation does not come at the cost of utilization.

All baselines experience high vLLM latency due to the lack of memory isolation. Similar to memory contention shown in Figure~\ref{fig:motiv:memory_study}, the colocated LP task can swap out vLLM memory during execution, leading to latency spikes. They also suffer from low LP throughput due to inefficient paging with the default GPU driver.

Co-locating vLLM and LlamaFactory (Figure~\ref{fig:eval:vllm+llama_factory_mempressure_bar}) also shows a similar trend. \tool reduces vLLM latency by 1.1-53$\times$, comparing against the best performing baseline.
\tool does not achieve the best LP throughput in this case, because it has to enforce a strict memory limit on LlamaFactory to accommodate bursty vLLM load.
Nevertheless, it still achieves the highest overall goodput.

Using the same \codeIn{gcgroup}-like interface, we also implement a dynamic policy, \tool{}DYN, which dynamically adjusts and relaxes the memory limit for the LP task to increase utilization, when vLLM has a low load that cannot fully saturate its reservation.
As Figure~\ref{fig:eval:vllm+diffusion_mempressure_bar} and Figure~\ref{fig:eval:vllm+llama_factory_mempressure_bar} show, \tool{}DYN slightly increases vLLM latency relative to \tool that enforces strict isolation, but still outperforms all baselines.
In return, \tool{}DYN improves LP throughput by an average of 2$\times$, indicating that modest slack in HP latency can be traded for substantial LP throughput gains.

% \yq{Trace: we construct bursty loads for LLM inference so vLLM will be idle periodically. \tool should be able to constrain the memory usage of the best-effort task when vLLM is running and harvest memory when vLLM is idle.}

% \yq{Setup: we will set a fixed memory size for HP and LP tasks. For example, 30GB for vLLM and 10GB for the other LP task.}

% \yq{Plot: we show three bar charts: vLLM p99 TTFT, vLLM P99 ITL, Best-effort throughput, similar to Figure~\ref{fig:motiv-mem-cont-perf}, where each bar represents a system. 
% We expect \tool to retain both much lower vLLM latency while achieve much higher offline throughput.}


\begin{figure*}[t]
    \centering
    \begin{subfigure}[b]{0.95\linewidth}
        \centering
        \includegraphics[width=\linewidth]{figs/vllm+diffusion_nomem_performance_bar.pdf}
        \vspace{-2em}
        \caption{Multimodal AI: vLLM (high priority) + Diffusion (low priority)}
        \label{fig:eval:vllm+diffusion_memsufficient_bar}
        \vspace{0.6em}
    \end{subfigure}
    
    \begin{subfigure}[b]{0.95\linewidth}
        \centering
        \includegraphics[width=\linewidth]{figs/vllm+llama_factory_nomem_performance_bar.pdf}
        \vspace{-2em}
        \caption{Continual online learning: vLLM (high priority) + LlamaFactory (low priority)}
        \label{fig:eval:vllm+llama_factory_memsufficient_bar}
    \end{subfigure}
    \vspace{-1em}
    \caption{Performance variation when co-running vLLM with (a) Diffusion and (b) LlamaFactory when memory is sufficient.}
    \label{fig:eval:memsufficient-overall}
    % \vspace{-1em}
\end{figure*}


\mysubsubsection{GPU Compute Efficiency}
\begin{comment}
We additionally evaluate whether \tool can achieve high GPU utilization and overall throughput even when GPU memory is sufficient in Figure~\ref{fig:eval:memsufficient-overall}. 
\yq{TODO after numbers are finalized. In text we will explain why GVM points are clustered into tree areas}
\end{comment}

We next isolate the impact of \tool's compute virtualization by running the same co-location experiments on an A100-80G GPU, so that memory contention is effectively removed.
We keep the other hardware and software setups the same to ensure this configuration exercises only compute-related mechanisms.
Results are shown in Figure~\ref{fig:eval:memsufficient-overall}.

% \hx{Where are the results?}
% Even with solely the scheduler being effective
% \joey{do you mean evaluated instead of effective}
In this case, \tool achieves the lowest latency for vLLM (the high-priority task) across both pipelines.
Specifically, \tool significantly outperforms MIG and achieves up to 480$\times$, 183$\times$, 2.75$\times$, and 2.28$\times$ lower latencies for median and P99 TTFT and ITL, respectively, across both pipelines. 
MIG performs worst in this case due to the NVIDIA hardware constraints that it can allocate up to 50\% of GPU compute and memory to vLLM, which is insufficient for vLLM to handle its peak load.
This highlights that \tool's driver-hardware collaborative scheduling provides much more flexible resource limit configurations than MIG and stronger isolation than the other baselines.
For LP task throughput, \tool{} performs slightly worse than MIG but outperforms all other baselines. Compared to MIG, \tool achieves 86.7\% Diffusion throughput for the multimodal AI workload and 88.3\% LlamaFactory throughput for continual learning workload. MIG achieves the best LP throughput because it allocates the rest 50\% GPU resources to the LP task.
% These results demonstrate that \tool{} achieves better overall utilization than MIG.
TGS achieves low HP latency but also much lower LP throughput because it conservatively stops scheduling LP kernels to provide compute isolation to vLLM. GPreempt and XSched, in contrast, aim for high overall throughput but fail to provide strong enough isolation to keep vLLM latency low.

Overall, \tool{} achieves an ideal balance between HP and LP tasks by keeping consistently low HP latency while significantly improving LP throughput.

% The reason \tool{} provides better compute isolation is that, unlike the baselines, it is the only system that supports true preemption.
% Other baselines must at least wait for the current timeslice of an application to finish, which delays their response to workload changes.
% In contrast, \tool{} can immediately preempt low-priority tasks, enabling much stronger isolation.
% \tool{} also achieves higher utilization because preemption in \tool{} happens at the GPU level, without requiring the CPU to maintain all scheduling information or to frequently communicate with the GPU.
% Avoiding such high-frequency CPU–GPU RPCs reduces overhead and leads to much better utilization.


\begin{figure*}[t]
    \centering
    \includegraphics[width=0.88\linewidth]{figs/vllm+diffusion_breakdown_performance.pdf}
    \vspace{-1.5em}
    \caption{All three of \tool's contributions work in tandem to improve latency and throughput.}
    \label{fig:eval:ablation}
    \vspace{-1em}
\end{figure*}




\mysubsection{Dynamic Resource Coordination}
\label{sec:eval:tradeoff}

In this experiment, we evaluate \tool's elasticity by dynamically adjusting resource limit and allocation between applications and investigate whether can constantly achieve Pareto frontier compared to baseline systems. 



Figure~\ref{fig:eval:vllm+diffusion-pareto} plots SLO attainment against normalized diffusion throughput when co-locating vLLM (HP) and Diffusion (LP).
We define the SLO target as vLLM’s P99 latency when running alone, and normalize Diffusion throughput to its ideal throughput on a dedicated GPU.

None of the baselines achieve more than 80\% SLO attainment or 0.35$\times$ normalized throughput.
This aligns with earlier results: they neither enforce memory isolation nor implement efficient scheduling and paging under contention.

\setlength{\columnsep}{8pt}
% \setlength{\wrapfigsep}{2pt}
\begin{wrapfigure}[15]{r}{0.45\linewidth}
    \centering
    \vspace{-0.8em}
    \includegraphics[width=\linewidth]{figs/vllm+diffusion_pareto.pdf}
    % \vspace{-2.5em}
    \caption{SLO attainment vs. throughput when co-locating vLLM (HP) and Diffusion (LP). \tool consistently stays on the Pareto frontier.}
    % \vspace{-1em}
    \label{fig:eval:vllm+diffusion-pareto}
\end{wrapfigure}

\tool, in contrast, achieves 100\% SLO attainment while matching the best LP throughput among baselines.
More importantly, \tool can adjust resource allocation according to HP latency requirements.
By varying the compute duty cycle of the LP task, \tool trades off vLLM latency against diffusion throughput and eventually reaches full LP throughput.
Across this range, \tool remains on the Pareto frontier, indicating that no other system dominates it on both metrics.

\tool{}DYN pushes this frontier further by dynamically reallocating GPU memory based on the current workload. It maintains over 90\% SLO attainment for high-priority requests while achieving an LP throughput that static \tool can reach only when its SLO attainment drops to around 20\%.
This demonstrates the effectiveness of coordinated adjustment of both compute and memory limits in improving utilization.

% \tool{}DYN pushes the Pareto frontier further outward.
% By dynamically allocating GPU memory based on the current workload, it effectively achieves our main goal: improving utilization.
% As shown in the figure, \tool{}DYN is able to maintain over 90\% SLO attainment for high-priority tasks, while reaching a level of low-priority throughput that the original \tool{} can achieve only when its high-priority SLO attainment drops to around 20\%.


\mysubsection{Performance Analysis}
\label{sec:eval:ablation}

We next investigate specific aspects of \tool's design to understand their individual contributions to overall performance.

\MyPara{Ablation Study.}
We incrementally enable \tool's mechanisms while co-locating vLLM and Diffusion under the same memory-constrained setting as in Figure~\ref{fig:eval:mempressure-overall}.
Figure~\ref{fig:eval:ablation} summarizes the results.

% Both isolation and efficient sharing mechanisms contribute greatly to the final performance. 
Adding memory isolation alone protects vLLM from heavy paging, substantially reducing vLLM tail latency by 29.5\% for P99 TTFT and 78.5\% for P99 ITL. 
Enabling compute isolation prioritizes vLLM's kernels, further reducing median and P99 TTFT by 99.4\% and 95.5\%, and median and P99 ITL by 52.5\% and 86.7\%, respectively.
Finally, enabling \tool's optimized paging increases Diffusion throughput by 69.4\%, which previously suffered from slow swapping. 
Interestingly, efficient paging also reduces vLLM latency. Our inspection reveals that the GPU cannot promptly switch process contexts until outstanding page faults are resolved, so a swap-heavy process can delay other processes even when the scheduler nominally assigns them proportional compute time. By reducing page-fault and swap time, \tool removes this hidden source of interference.

\MyPara{Microbenchmarks.}
We further quantify the impact of \tool's swapping optimizations on low-priority tasks under memory contention.
Figure~\ref{fig:eval:gvm-swap-breakdown} shows the time breakdown of diffusion when its available GPU memory is capped to 40\% of its nominal requirement.
With \tool, memory swapping overhead is comparable to compute time, resulting in a manageable 2$\times$ throughput reduction.
In contrast, the unoptimized baseline suffers a severe drop ($>$15$\times$ in Figure \ref{fig:motiv:swap_breakdown}) due to synchronous blocking and small-packet PCIe congestion.
% In other words, memory-swapping causes an overall 2$\times$ reduction in throughput.
% In contrast, the unoptimized memory-swapping baseline (Figure \ref{fig:motiv:swap_breakdown}) suffers a 15$\times$+ drop in throughput under the same memory limit. \yc{Looks suspicious because Figure \ref{fig:eval:diffusion-mem-tput} doesn't show such a big difference in any given configuration}.
% 
Figure~\ref{fig:eval:diffusion-mem-tput} plots diffusion throughput versus its GPU memory usage.
Notably, \tool consistently outperforms the unoptimized baseline across all configurations, and can offload 10\% of its working set with minor throughput drop.

% With our optimizations, allocating just 90\% of the required memory is sufficient to achieve 100\% of the original performance, demonstrating the effectiveness of overlapping compute with memory transfer.


\begin{figure}[t]
    \centering
    \begin{minipage}[t]{0.46\linewidth}
        \vspace{0pt}
        \centering
        \includegraphics[width=\linewidth]{figs/gvm_swap_breakdown.pdf}
        \vspace{-2.2em}
        \caption{Time breakdown of Diffusion under memory pressure.}
        \label{fig:eval:gvm-swap-breakdown}
        \vspace{-0.8em}
    \end{minipage}
    \hfill
    \begin{minipage}[t]{0.49\linewidth}
        \vspace{0pt}
        \centering
        \includegraphics[width=\linewidth]{figs/swap_optimization.pdf}
        \vspace{-2.2em}
        \caption{Normalized Diffusion throughput with regard to GPU memory usage.}
        % \vspace{-0.8em}
        \label{fig:eval:diffusion-mem-tput}
    \end{minipage}
\end{figure}

